\section{Consequences for closed systems and semidefinite relaxations}
\label{sec:consequences}

The certificate theorem decides whether the hull is the whole space.
Its consequences below concern this dichotomy; they do not identify
every supporting inequality of a proper hull.

\subsection{Closed inequalities and the three feasibility regimes}

Write $T=\{x:f_i(x)\leq0\text{ for all }i\}$.

\begin{corollary}\label{cor:closed}
Suppose $S\ne\varnothing$ and $f^h$ has asymptotic hyperplane convexity.
Then
\[
 \conv T\ne\R^n\quad\Longleftrightarrow\quad
 \conv S\ne\R^n\quad\Longleftrightarrow\quad
 \mathcal K\not\subseteq\mathcal T.
\]
\end{corollary}
\begin{proof}
The inclusion $S\subseteq T$ gives one implication. A nontrivial convex
certificate bounds $T$ by the closed convex set $\{f_\lambda\leq0\}$.
That set is proper because the nonconstant convex quadratic
$f_\lambda$ is unbounded above. The remaining implication is
Theorem~\ref{thm:certificate}.
\end{proof}

Strict feasibility is essential here; Example~\ref{ex:closed} gives a
counterexample even with HHC and a closed-system hull with interior.
No equality between $\conv T$ and $\cl\conv S$ is asserted.

\begin{corollary}[Three regimes under HHC]\label{cor:regimes}
Suppose $f^h$ has HHC and $n\geq2$. Exactly one of the following holds:
\begin{enumerate}
\item $S=\varnothing$. This is equivalent to the existence of
      $\lambda\geq0$, $\lambda\ne0$, with $Q_\lambda\succeq0$.
\item $S\ne\varnothing$ and $\conv S=\R^n$. This is equivalent to
      nonemptiness together with $\mathcal K\subseteq\mathcal T$.
\item $S\ne\varnothing$ and $\conv S\ne\R^n$. This is equivalent to
      nonemptiness together with $\mathcal K\not\subseteq\mathcal T$.
      If also $n\geq3$, the good aggregations defined in
      Section~\ref{sec:setting} describe $\conv S$ exactly.
\end{enumerate}
\end{corollary}
\begin{proof}
For the first statement, HHC implies convexity of $f^h(\R^{n+1})$
because $n+1\geq3$. Moreover, $S=\varnothing$ if and only if
$S^h=\varnothing$: a strict homogeneous point with $t=0$ can be
perturbed to one with $t\ne0$, and then dehomogenized.
Separation of the convex homogeneous image from the open negative
orthant gives a nonzero $\lambda\geq0$ for which
$\lambda\transpose f^h\geq0$, equivalently $Q_\lambda\succeq0$.
The converse follows by evaluation at a hypothetical strict point.
This is the classical separation argument used in
\citet[Proposition~2.13]{BDSpreprint}.
Theorem~\ref{thm:certificate} supplies the second and third statements.
The final hull description is
\citet[Theorem~2.9]{BDSpreprint}.
\end{proof}

The proper-hull certificate itself remains valid for $n=1$ by
Theorem~\ref{thm:certificate}; the dimension restrictions above belong
to the additional emptiness and full-description statements.

\subsection{A finite semidefinite characterization}

The certificate cone can be examined on the compact slice
\[
 \mathcal D=\{\lambda\geq0:\textstyle\sum_i\lambda_i=1,
                                      \ A_\lambda\succeq0\}.
\]
If $\mathcal D$ is nonempty, let
\begin{equation}\label{eq:sdp-objectives}
 \tau=\max_{\lambda\in\mathcal D}\tr A_\lambda,
 \qquad
 \beta_j^\pm=\max_{\lambda\in\mathcal D}\bigl(\pm(b_\lambda)_j\bigr)
 \quad(j=1,\ldots,n).
\end{equation}
All maxima are attained. These are $2n+1$ semidefinite programs, with
one $n\times n$ positive semidefinite block and $m$ nonnegative weights.

\begin{corollary}\label{cor:finite-sdp}
The inclusion $\mathcal K\subseteq\mathcal T$ holds exactly when
$\mathcal D=\varnothing$, or $\mathcal D\ne\varnothing$ and all the
values in~\eqref{eq:sdp-objectives} are zero. Under the hypotheses of
Theorem~\ref{thm:certificate}, these conditions characterize
$\conv S=\R^n$.
\end{corollary}
\begin{proof}
Every nonzero element of $\mathcal K$ has a positive multiple in
$\mathcal D$. For a positive semidefinite matrix, zero trace is
equivalent to the matrix being zero. Thus $\tau=0$ forces
$A_\lambda=0$ throughout $\mathcal D$. The two coordinate objectives
force $(b_\lambda)_j\leq0$ and $(b_\lambda)_j\geq0$ there. The converse
is immediate. If $\mathcal D$ is empty, $\mathcal K=\{0\}$.
\end{proof}

The reduction uses the trace in place of all matrix coordinates.
Feasibility of $\mathcal D$ is the common feasibility part of these
programs. The characterization uses exact feasibility and optimal values;
it is not a bit-complexity result or a floating-point decision procedure,
and small computed objective values do not establish the exact zeros
needed for a whole-space conclusion.

\subsection{What the Shor projection preserves}

Define the projection of the standard semidefinite relaxation by
\begin{equation}\label{eq:shor}
 P_{\rm SDP}=\left\{x\in\R^n:\begin{array}{l}
       \exists X\in\Sym^n,\quad X-xx\transpose\succeq0,\\
       \langle A_i,X\rangle+2b_i\transpose x+c_i\leq0
                      \quad(i=1,\ldots,m)
      \end{array}\right\}.
\end{equation}
The matrix condition is equivalent to
$\left(\begin{smallmatrix}1&x\transpose\\x&X\end{smallmatrix}\right)
\succeq0$, so $P_{\rm SDP}$ is convex. It contains $T$, by taking $X=xx\transpose$.
Let
\begin{equation}\label{eq:convex-aggregation-intersection}
 C_{\rm agg}=\bigcap_{\lambda\in\mathcal K}\{x:f_\lambda(x)\leq0\}.
\end{equation}
This is a closed convex set. For any Shor witness $X$ and
$\lambda\in\mathcal K$,
\begin{equation}\label{eq:shor-aggregation}
 f_\lambda(x)=
 \sum_i\lambda_i\bigl(\langle A_i,X\rangle+2b_i\transpose x+c_i\bigr)
       -\langle A_\lambda,X-xx\transpose\rangle\leq0.
\end{equation}
Consequently $P_{\rm SDP}\subseteq C_{\rm agg}$, the classical
relationship between semidefinite relaxation and convex aggregation
studied by \citet{FujieKojima1997,KojimaTuncel2000}. We prove the
closure statement needed here without assuming that the projection is
closed.

\begin{proposition}\label{prop:shor-closure}
If $S\ne\varnothing$, then $\cl P_{\rm SDP}=C_{\rm agg}$. Under the same
assumption, without any hidden-convexity hypothesis,
\begin{equation}\label{eq:shor-whole}
 P_{\rm SDP}=\R^n\quad\Longleftrightarrow\quad
 \mathcal K\subseteq\mathcal T.
\end{equation}
\end{proposition}
\begin{proof}
Put $g(v)=(v\transpose A_i v)_i$ and introduce the convex cone
\[
 \mathcal C=\{(\langle A_i,Y\rangle)_i:Y\succeq0\}+\R^m_+.
\]
It contains $\R^m_+$ and is therefore full dimensional. It need not
be closed. Directly from~\eqref{eq:shor},
\[
 x\in P_{\rm SDP}\ \Longleftrightarrow\ -f(x)\in\mathcal C,
 \qquad
 \mathcal C^*=\mathcal K,
 \qquad
 x\in C_{\rm agg}\ \Longleftrightarrow\ -f(x)\in\cl\mathcal C.
\]
The dual identity follows from the self-duality of the positive
semidefinite cone; the last equivalence is the finite-dimensional
bipolar theorem.

We will use
$\cl\mathcal C+\interior\mathcal C\subseteq\interior\mathcal C$.
For completeness, if $B(c,r)\subseteq\mathcal C$ and
$d_k\in\mathcal C$ tends to $d\in\cl\mathcal C$, then for large $k$
the ball $B(c+d,r/2)$ lies in $B(c+d_k,r)$, which lies in
$\mathcal C+\mathcal C=\mathcal C$. This also proves that adding
another element of $\mathcal C$ preserves the interior.

Choose $x_0\in S$, so $-f(x_0)\in\R^m_{++}\subseteq
\interior\mathcal C$. For $x\in C_{\rm agg}$ and $0<t<1$, put
$x_t=(1-t)x+t x_0$. The quadratic identity
\begin{equation}\label{eq:quadratic-mixing}
 -f(x_t)=(1-t)(-f(x))+t(-f(x_0))
                         +t(1-t)g(x-x_0)
\end{equation}
puts $-f(x_t)$ in $\interior\mathcal C$, since
$g(v)=(\langle A_i,vv\transpose\rangle)_i\in\mathcal C$.
Thus $x_t\in P_{\rm SDP}$, and $x_t\to x$ proves
$C_{\rm agg}\subseteq\cl P_{\rm SDP}$. The opposite inclusion follows from
\eqref{eq:shor-aggregation} and closedness of $C_{\rm agg}$.

A nontrivial convex certificate makes $C_{\rm agg}$ proper, so
$P_{\rm SDP}=\R^n$ implies $\mathcal K\subseteq\mathcal T$.
Conversely, suppose every certificate is trivial. For any
$\lambda\in\mathcal K\setminus\{0\}$ strict feasibility gives
$f_\lambda\equiv c_\lambda<0$. Hence $C_{\rm agg}=\R^n$ and
$-f(z)\in\cl\mathcal C$ for every $z$.
For an arbitrary target $x$, apply~\eqref{eq:quadratic-mixing} with
$z=2x-x_0$ in place of $x$ and $t=1/2$. Its left side is $-f(x)$;
its right side belongs to $\interior\mathcal C$. Thus $x\in P_{\rm SDP}$.
\end{proof}

The closure conclusion also follows from
\citet[Theorem~2.1]{FujieKojima1997}: strict feasibility provides
their semidefinite Slater condition by choosing
$X=x_0x_0\transpose+\varepsilon I$ for sufficiently small
$\varepsilon>0$. The cone proof above shows where strict feasibility
enters and why no cone image needs to be closed.

\begin{corollary}\label{cor:shor-hull}
If $S\ne\varnothing$ and $f^h$ has asymptotic hyperplane convexity,
then
\[
 P_{\rm SDP}=\R^n\quad\Longleftrightarrow\quad\conv S=\R^n
             \quad\Longleftrightarrow\quad\conv T=\R^n.
\]
\end{corollary}
\begin{proof}
Combine Proposition~\ref{prop:shor-closure},
Theorem~\ref{thm:certificate}, and Corollary~\ref{cor:closed}.
\end{proof}

Under these hypotheses the relaxation thus preserves the existence of
a nonconstant valid affine bound, but not necessarily each particular
bound (Example~\ref{ex:strip}).

\begin{example}[A nonclosed Shor projection with strict feasibility]
\label{ex:nonclosed-shor}
In $\R^2$, let
\[
 f_1(x)=2x_1x_2+1,\qquad f_2(x)=x_1^2-x_1.
\]
The point $(1/2,-2)$ is strictly feasible. A nonnegative aggregate has
quadratic matrix
$\left(\begin{smallmatrix}\lambda_2&\lambda_1\\
\lambda_1&0\end{smallmatrix}\right)$, which is positive semidefinite
exactly when $\lambda_1=0$. Therefore
\[
 C_{\rm agg}=[0,1]\times\R,
 \qquad
 P_{\rm SDP}=((0,1)\times\R)\ \cup\ (\{1\}\times(-\infty,-1/2]).
\]
To verify the latter equality, write $Y=X-xx\transpose\succeq0$.
The second lifted row gives $0\leq Y_{11}\leq x_1-x_1^2$, so
$0\leq x_1\leq1$. At either endpoint $Y_{11}=0$, and positive
semidefiniteness forces $Y_{12}=0$. At $x_1=0$ the first row then
fails; at $x_1=1$ it requires $x_2\leq-1/2$, and $Y=0$ suffices.
For $0<x_1<1$, take
\[
 X_{11}=\frac{x_1^2+x_1}{2},\quad X_{12}=-\frac12,
 \quad
 Y_{22}=\frac{(-1/2-x_1x_2)^2}{(x_1-x_1^2)/2},
 \quad X_{22}=x_2^2+Y_{22}.
\]
Then $Y\succeq0$ has zero determinant and the lifted rows hold.
The displayed description proves nonclosedness even with strict
feasibility and, by Dines' theorem, HHC.

Compactness of the original closed feasible set does not in general
restore projection closedness. Add the two rows
$f_3=-2x_1x_2-2$ and $f_4=1/4-x_1$ to this example.
The point $(1/2,-3/2)$ is strictly feasible and
\[
 T=\{x:1/4\leq x_1\leq1,\ -1\leq x_1x_2\leq-1/2\}
\]
is closed and bounded. A nonnegative aggregate is convex exactly
when $\lambda_1=\lambda_3$, and then equals
\[
 \lambda_2(x_1^2-x_1)+\lambda_4(1/4-x_1)-\lambda_1.
\]
It follows that $C_{\rm agg}=[1/4,1]\times\R$. In contrast,
\[
 P_{\rm SDP}=([1/4,1)\times\R)\ \cup\ (\{1\}\times[-1,-1/2]).
\]
Indeed the linear fourth row requires $x_1\geq1/4$;
for $x_1<1$ the preceding lift works with $X_{12}=-3/4$ and
$Y_{22}=(-3/4-x_1x_2)^2/((x_1-x_1^2)/2)$.
At $x_1=1$, $Y_{11}=Y_{12}=0$, so the first and third rows require
$-1\leq x_2\leq-1/2$, and a rank-one lift suffices.
This compact variant only shows that a closure qualification is
necessary; no HHC claim is made for it.
\end{example}

We use the classical SDP/aggregation relationship only through the
closure equality in Proposition~\ref{prop:shor-closure} and
\citet[Theorem~2.1]{FujieKojima1997}, not through the unqualified
projection equality stated in \citet[Theorem~4.2]{KojimaTuncel2000}.
Example~\ref{ex:nonclosed-shor} explains the distinction, including
when the original feasible set is compact. For that compact variant, the cone generated by its four defining
quadratic functions is a closed convex cone by finite generation
(Lemma~\ref{lem:finite-cone}). Replacing the finite list by this cone
leaves the original feasible set, the semidefinite projection, and the
convex-aggregation intersection unchanged: the original rows belong to
the cone, and every cone member is a nonnegative combination of them.
Thus it also meets the closed-convex-cone representation hypothesis in
that theorem.
